\documentclass[10pt,a4paper]{amsart}
\usepackage[margin=19mm]{geometry}
\usepackage{amsmath,amssymb,mathtools}
\usepackage[scr=rsfs,cal=euler]{mathalfa}
\usepackage{xfrac}
\usepackage[unicode,hidelinks,bookmarks=false]{hyperref}
\newcommand{\ie}{i.e.\ }
\newcommand{\cf}{cf.\ }
\newcommand{\resp}{respectively\ }
\newcommand{\Subsection}{\subsection}
\providecommand{\nobreakdash}{\nobreak\hbox{-}}
\setlength{\emergencystretch}{2em}
\multlinegap0pt
\usepackage{bm}
\usepackage{bbm}
\usepackage{dsfont}
\usepackage{xspace}
\usepackage{cancel}
\usepackage{mathtools}
\usepackage{amsthm}
\makeatletter
\@ifundefined{thm}{\newtheorem{thm}{Theorem}}{}
\@ifundefined{lem}{\newtheorem{lem}[thm]{Lemma}}{}
\makeatother
\title[On the Sum of Additive Characters and its Applications over ]{On the Sum of Additive Characters and its Applications over Finite Fields}
\author{Maithri K., Vadiraja Bhatta G. R.,, Indira K. P}
\date{Source: arXiv:2506.16081v2 (CC BY 4.0)}
\begin{document}
\maketitle

\noindent\emph{Source and licence.} On the Sum of Additive Characters and its Applications over Finite Fields. Version arXiv:2506.16081v2 (CC BY 4.0), distributed under the Creative Commons Attribution 4.0 International licence, as recorded in the arXiv licence metadata for this exact version and shown on the abstract page https://arxiv.org/abs/2506.16081v2. Full rights notice: licenses/arxiv-2506.16081-CC-BY-4.0.txt.\par\smallskip

\noindent\textit{Editorial scope.} Counted unit: the Thm whose source label is \texttt{MainThm} in the article (source0.tex lines 214--217, source label \texttt{MainThm}), delivered together with the article's own complete original proof of it (source0.tex lines 218--279, one \texttt{proof} environment of 62 lines). No mathematics is written, repaired or re-derived: statement and proof are the article's own text line for line. Required context retained uncounted: Sumaddi (lines 114--117); Lemmasumphi (lines 173--175); Lemmaphi(hg) (lines 184--186). Every definition, display and label of those lines is retained unchanged. Omitted from this package and not redistributed: 1--113, 118--172, 176--183, 187--213, 280--397. Nothing inside a retained excerpt is omitted or altered: every retained body is delivered exactly as the article's lines are written, so no omission receipt of any kind accompanies this package. Citations: the retained excerpts carry no citation command at all, so no bibliography entry of the article is transcribed and no reference is redistributed. Reference adaptation: none was needed and none was made; every reference inside the delivered unit resolves inside it, so no unresolved reference or citation is produced. Printed slips kept literal: no printed word, symbol, hypothesis or number of the counted statement or of its proof was altered. Layout and class: the article's own class is \texttt{article} with packages including lmodern, xcolor, amsthm, amsmath, amsfonts, amssymb, anysize, array, tabularx; the wrapper is the lane's pinned \texttt{amsart} editorial wrapper at 10pt on A4, which restates the theorem-like environments the retained text needs and adds the article's own macro definitions that the retained text needs (none), with extra wrapper packages (bm, bbm, dsfont, xspace, cancel, mathtools). The delivered numbering is the unit's own. Assets: no retained span contains a figure environment, a graphics-inclusion command, a PDF-page inclusion, a table or a TikZ picture; no image file is copied into this package, the package declares no asset dependency and no image file is redistributed with it. Licence: version arXiv:2506.16081v2 of this article is distributed under the Creative Commons Attribution 4.0 International licence, as recorded in the arXiv licence metadata for this exact version and shown on the abstract page https://arxiv.org/abs/2506.16081v2. A full rights notice is delivered at licenses/arxiv-2506.16081-CC-BY-4.0.txt. Characters: every retained excerpt is pure ASCII, so the delivered engine, XeLaTeX with its default Latin Modern text font, reports no missing character.
    \begin{thm}\label{Sumaddi}
        Let $x^m-1 = f_1(x)f_2(x), ~g(x) \vert x^m-1$ and let $\alpha \in \mathbb{F}_{q^m}$ with $\mathbb{F}_q$-Order $f_1$, then
        $\displaystyle\sum_{g} \chi (\alpha) = \mu \left (d \right)\frac{\phi(g)}{\phi\left(d\right)}$, where $d(x) = \dfrac{g(x)}{gcd \left(g(x),f_2\right)}$ and the summation runs over all the additive characters of $\mathbb{F}_q$-Order $g$. 
    \end{thm}

    \begin{lem}\label{Lemmasumphi}
         If $u$ is an irreducible polynomial over $\mathbb{F}_q[x]$ of degree $n$, then $\displaystyle \sum_{i=0}^l\phi(u^i) = \frac{\phi(u^{l+1})}{\phi(u)}$ for any positive integer $l$.
     \end{lem}

    \begin{lem}\label{Lemmaphi(hg)}
        Let $f(x) \in \mathbb{F}_q[x]$ be a divisor of $x^m-1$. For any square-free divisor $h(x)$ of $f, \displaystyle \sum_{\substack{g(x) \vert \frac{x^m-1}{f}\\ gcd(h,\frac{x^m-1}{fg}) =1} } \phi(hg) = q^{deg(\frac{x^m-1}{f})}\phi(h)$.
    \end{lem}

     \begin{thm}\label{MainThm}
         Let $f(x)$ be a divisor of $x^m-1$ of degree $n$ over $\mathbb{F}_q[x]$. Then, for any $\alpha \in \mathbb{F}_{q^m}$,
         $$\eta_f(\alpha) = \displaystyle \frac{\phi(f)}{q^m}\sum_{h\vert f} \frac{\mu(h)}{\phi(h)} \sum_{\substack{g \vert \frac{x^m-1}{f}\\ gcd(h,\frac{x^m-1}{fg}) =1} } \sum_{hg} \chi(\alpha) $$  where the inner sum is running over all the additive characters of $\mathbb{F}_q$-Order $hg$.
     \end{thm}

     \begin{proof}
        Let $k(x)$ be the $\mathbb{F}_q$-Order of $\alpha$ and $x^m-1=(p_1(x)p_2(x)\dots p_r(x))^a, f=p_1^{b_1} p_2^{b_2}\dots p_r^{b_r}, k=p_1^{c_1} p_2^{c_2}\dots p_r^{c_r}$. 

        By using theorem \ref{Sumaddi}, 
        \begin{equation}\label{EQThm}
            \eta_f(\alpha) = \displaystyle \frac{\phi(f)}{q^m}\sum_{h\vert f} \frac{\mu(h)}{\phi(h)} \sum_{\substack{g \vert \frac{x^m-1}{f}\\ gcd(h,\frac{x^m-1}{fg}) =1} } \mu \left( \frac{gh}{\text{gcd}(gh,\frac{x^m-1}{k})}\right)\frac{\phi(gh)}{\phi\left( \frac{gh}{\text{gcd}(gh,\frac{x^m-1}{k})}\right)}
        \end{equation}
         
        The following three mutually exclusive cases cover all possibilities:\\
         \textbf{Case 1:} $k(x)=f(x)$\\
         
         Since gcd$(h,\frac{x^m-1}{fg}) = 1$, gcd$(hg,\frac{x^m-1}{f})=g$. The equation \ref{EQThm} reduces to,
         $$\eta_f(\alpha)= \frac{\phi(f)}{q^m}\sum_{h\vert f} \frac{\mu(h)}{\phi(h)} \sum_{\substack{g \vert \frac{x^m-1}{f}\\ gcd(h,\frac{x^m-1}{fg}) =1} } \mu(h) \frac{\phi(gh)}{\phi(h)} $$
         By using lemma \ref{Lemmaphi(hg)}, the above expression becomes
         \begin{align*}
             \eta_f(\alpha) &=\frac{\phi(f)}{q^m}\sum_{h\vert f} \frac{(\mu(h))^2}{(\phi(h))^2}q^{m-n}\phi(h)=\frac{\phi(f)}{q^m}q^{m-n}\sum_{h\vert f} \frac{(\mu(h))^2}{\phi(h)} \\
             & =\frac{\phi(f)}{q^m}q^{m-n}\frac{q^n}{\phi(f)} = 1.
         \end{align*}
         \textbf{Case 2:} $k \mid f$ and $k \neq f$\\
         Without loss of generality we assume that $f \neq 1$, we take $b_r \neq 0$.\\
        Let $S_h = \displaystyle\frac{\mu(h)}{\phi(h)}\sum_{\substack{g \vert \frac{x^m-1}{f} \\ \text{gcd}(h,\frac{x^m-1}{fg})=1}} \mu\left( \frac{gh}{\text{gcd}(gh,\frac{x^m-1}{k})}\right)\frac{\phi(gh)}{\phi\left( \frac{gh}{\text{gcd}(gh,\frac{x^m-1}{k}}\right)}$.\\
        Suppose, $c_i<b_i ~ \forall i, 0 \leq i \leq r.$
        Since $g \vert \frac{x^m-1}{f}, c_i<b_i$ and $h$ is a square-free divisor of $f$, it follows that $gh \vert \frac{x^m-1}{k}$. Thus, gcd$(gh,\frac{x^m-1}{k})= gh.$\\
        Then $S_h = \dfrac{\mu(h)}{\phi(h)} \displaystyle \sum_{\substack{g \vert \frac{x^m-1}{f} \\ \text{gcd}(h,\frac{x^m-1}{fg})=1}} \phi(gh) = \mu(h) q^{m-n}$.\\
        Now for any divisor $h$ of $f$, let us define $h^c = \begin{cases}
            \frac{h}{p_r} \text{ if } p_r \vert h \\
            hp_r \text{ otherwise}.
        \end{cases}$\\
        Note that both $h$ and $h^c$ are square-free divisors of $f$, and they are distinct. Moreover, we have $\mu(h^c) = -\mu(h),$
        thus $S_h + S_{h^c} = q^{m-n}(mu(h) + mu(h^c)) = 0$.\\
        Thus, all such terms cancel in pairs, and the total contribution to $\eta_f(\alpha)$ is zero.\\
        Suppose, $c_1=b_1, c_2 = b_2, \dots, c_j=b_j$ for some $j,~1\leq j \leq r$. Note that at least one $c_j \neq b_j$. Let us take $c_r \neq b_r$. \\
        Now, gcd$\left( gh , \dfrac{x^m-1}{k}\right) = \dfrac{gh}{\text{gcd}(h,p_1p_2\dots p_j)}$. Then,
        \begin{align*}
            S_h &= \frac{\mu(h)}{\phi(h)} \sum_{\substack{g\vert \frac{x^m-1}{f}\\\text{gcd}(h,\frac{x^m-1}{fg})=1}} \frac{\mu(\text{gcd}(h,p_1p_2\dots p_j))}{\phi(\text{gcd}(h,p_1p_2\dots p_j))} \phi(hg) \\
            &= \frac{\mu(h)\mu(\text{gcd}(h,p_1p_2\dots p_j))}{\phi(\text{gcd}(h,p_1p_2\dots p_j))}q^{m-n}.
        \end{align*}
        Again, we pair each $S_h$ with a unique $S_{h^c}$, using the same rule for $h^c$ as above (adding/removing the fixed prime $p_r$). This time, note that
        \[
        \gcd(h^c, p_1p_2\dots p_k) = \gcd(h, p_1p_2\dots p_k),
        \]
        so the pre-factor is the same in both $S_h$ and $S_{h^c}$, and again\\
        $S_h + S_{h^c} = q^{m-n} \dfrac{\mu(\text{gcd}(h,p_1p_2\dots p_k))}{\phi(\text{gcd}(h,p_1p_2\dots p_k))} (\mu(h)+\mu(h^c)) = 0$.\\
        Thus, $\eta_f(\alpha) = 0$.\\
        \textbf{Case 3:} $k \nmid f$\\
        Then there exists an irreducible polynomial $u(x)$ such that $l$ is the highest power of $u$ dividing $k$ and $u^l$ does not divide $f$.\\
        Let $a$ be the highest power of $u$ dividing $x^m-1$, $r$ be the highest power of $u$ dividing $f$ and $h, g$ be fixed with $u \nmid g$, $h$  square-free. Let,
        $$S_g = \mu\left(\dfrac{gh}{\text{gcd}(gh,\frac{x^m-1}{k})}\right)\dfrac{\phi(gh)}{\phi\left(\dfrac{gh}{\text{gcd}(gh,\frac{x^m-1}{k})}\right)}.$$
        Consider the sum, $S = \displaystyle \sum_{i=0}^{a-r}S_{u^ig} = \sum_{i=0}^{a-r} \mu\left(\frac{u^igh}{\text{gcd}(u^igh,\frac{x^m-1}{k})}\right)\frac{\phi(u^igh)}{\phi\left(\frac{u^igh}{\text{gcd}(u^igh,\frac{x^m-1}{k})}\right)}$. \\
        Suppose $u\nmid h$ then,  $\text{gcd}(u^igh,\frac{x^m-1}{k}) = u^i \text{gcd}(gh,\frac{x^m-1}{k})$ for $0 \leq i \leq a-l$ and $\text{gcd}(u^igh,\frac{x^m-1}{k}) = u^{a-l} \text{gcd}(gh,\frac{x^m-1}{k})$ for $a-l < i \leq a-r$. Thus,
        
        \begin{align*}
            S =& \displaystyle \sum_{i=0}^{a-l} \mu \left( \frac{gh}{\text{gcd}(gh,\frac{x^m-1}{k})}\right) \frac{\phi(u^i)\phi(gh)}{\phi\left( \frac{gh}{\text{gcd}(gh,\frac{x^m-1}{k})}\right)} \\&+ \sum_{j=a-l+1}^{l} \mu(u^{j-a+l})  \mu \left( \frac{gh}{\text{gcd}(gh,\frac{x^m-1}{k})}\right) \frac{\phi(u^j)\phi(gh)}{\phi(u^{j-a+l})\phi\left( \frac{gh}{\text{gcd}(gh,\frac{x^m-1}{k})}\right)}   \\
            &= \mu \left( \frac{gh}{\text{gcd}(gh,\frac{x^m-1}{k})}\right) \frac{\phi(gh)}{\phi\left( \frac{gh}{\text{gcd}(gh,\frac{x^m-1}{k})}\right)}\left[1+\phi(u)+\phi(u^2)+\dots+ \phi(u^{a-l})-\frac{\phi(u^{a-l+1})}{\phi(u)}\right]
        \end{align*}
        By lemma \ref{Lemmasumphi}, we get $S=0$.\\
        
        Suppose $u \mid h$, as gcd$(h,\frac{x^m-1}{fg}) = 1, u^{a-r}$ must divide $g$, a contradiction to the choice of $g$ such that $u \nmid g$.\\           
        Thus, $\eta_f = 0$ .\\
        Similarly, if $u \mid g$, write $g = u^j t$ with $\gcd(t,u)=1$.  Replacing $g$ by $t$ in the above argument shows again that all powers of $u$ would have to divide $t$, contradicting $u\nmid t$.  Therefore this case also yields $\eta_f = 0$.

     \end{proof}

\end{document}
